\section{Lexical Semantics, WSD and NER}
\label{sec:lexical-semantics}
\slidesrc{Semantics (long version): 11--43}

\begin{nbwarning}[source of this section]
This section comes from the \emph{long} version of the semantics slides, which was not
shown in class (the lessons used the shortened decks of
Sections~\ref{sec:distributional-semantics} and~\ref{sec:frame-model-semantics}). It is
included because it completes the semantics part of the course map
(Section~\ref{sec:course-map}): lexical semantics, ontologies, word sense
disambiguation and named entity recognition. Check with the lecturer how much of it is
exam material.
\end{nbwarning}

Distributional semantics learns meaning from usage; frame and model-theoretic semantics
build structured meanings of sentences. \textbf{Lexical semantics} studies the meanings
of \emph{words} and the relations among them, and organises them into resources such as
WordNet. Two classic tasks use these resources: deciding which meaning of an ambiguous
word is meant (WSD), and finding the names of people, places and organisations in a text
(NER).

This section covers:
\begin{enumerate}[nosep]
  \item the lexicon and the relations among word meanings
        (Section~\ref{sec:lexical-relations});
  \item thesauri and WordNet (Section~\ref{sec:wordnet});
  \item word sense disambiguation and the Lesk algorithm
        (Section~\ref{sec:word-sense-disambiguation});
  \item named entity recognition (Section~\ref{sec:named-entity-recognition}).
\end{enumerate}

\subsection{The Lexicon and Lexical Relations}
\label{sec:lexical-relations}
\slidesrc{Semantics (long version): 11--21}

\subsubsection{Lexicon and lexical units}
\slidesrc{Semantics (long version): 11--13}

Linguistic theories generally regard a human language as made of two parts:
\begin{itemize}[nosep]
  \item a \textbf{lexicon}: essentially a catalogue of the language's words (its
        \emph{wordstock}). It is also thought to include bound morphemes, which cannot
        stand alone as words (most affixes, Section~\ref{sec:morphology}), and in some
        analyses compound words and certain idioms and collocations;
  \item a \textbf{grammar}: a system of rules that combine those words into meaningful
        sentences (Section~\ref{sec:grammars-parsing}).
\end{itemize}
Dictionaries try to list the lexicon in alphabetical order, usually without the bound
morphemes.

\begin{nbdef}{Lexical semantics}
The units of analysis of \textbf{lexical semantics} are \textbf{lexical units}: not only
words but also sub-words (affixes) and larger units (compound words, phrases). Lexical
units make up the lexicon. Lexical semantics studies how the meaning of lexical units
correlates with the structure of the language (syntax), and how meanings relate to each
other through \textbf{lexical relations}: hyponymy, hypernymy, synonymy, antonymy,
homonymy, \ldots
\end{nbdef}

\subsubsection{Synonymy}
\slidesrc{Semantics (long version): 14}

\begin{nbdef}{Synonymy}
\textbf{Synonymy} is similarity of meaning. Leibniz: two expressions are synonymous if
substituting one for the other \emph{never} changes the truth value of a sentence in which
the substitution is made.
\end{nbdef}

Such \emph{global} synonymy is rare (a language with perfect synonyms would be
redundant). A more useful notion is \textbf{synonymy relative to a context}: two
expressions are synonymous in a linguistic context $C$ if substituting one for the other
in $C$ does not alter the truth value. E.g.\ \emph{board} and \emph{plank} are
synonyms in ``he nailed a \underline{\quad} to the wall'' but not in ``the
\underline{\quad} of directors''. A consequence of defining synonymy by substitution:
words of different syntactic categories cannot be synonyms (\emph{quick} and
\emph{quickly} cannot replace each other).

\subsubsection{Antonymy}
\slidesrc{Semantics (long version): 15}

The antonym of a word $x$ is sometimes \emph{not-$x$}, but not always: \emph{rich} and
\emph{poor} are antonyms, but \emph{not rich} does not imply \emph{poor} (many people
consider themselves neither). Antonymy is a lexical relation between \emph{word forms},
not a semantic relation between concepts: \emph{rise}/\emph{fall} and
\emph{ascend}/\emph{descend} are antonym pairs, but \emph{rise}/\emph{descend} sounds
odd, even though the concepts are the same.

\subsubsection{Hyponymy and hypernymy}
\slidesrc{Semantics (long version): 16--17}

\begin{nbdef}{Hyponymy and hypernymy}
\textbf{Hyponymy} is a semantic relation between word meanings: \{maple\} is a
\textbf{hyponym} of \{tree\}. The inverse relation is \textbf{hypernymy}: \{tree\} is a
\textbf{hypernym} of \{maple\}. Also called subordination/superordination,
subset/superset, or the \textbf{IS-A} relation. Test: a native speaker must accept
sentences of the form ``An $x$ is a (kind of) $y$''. For verbs the relation is called
\textbf{troponymy} (\emph{to lisp} is a manner of \emph{to talk}).
\end{nbdef}

\begin{center}
\begin{tikzpicture}[w/.style={font=\small, text=nbblue}, e/.style={draw=nborange!80}]
  \node[w] (an) at (-1.2,0) {animals};
  \node[w] (do) at (-3.8,-1.1) {dogs};
  \node[w] (ca) at (0.4,-1.1) {cats};
  \node[w] (ho) at (2.0,-1.1) {horses};
  \node[w] (pi) at (3.5,-1.1) {pigs};
  \node[w] (te) at (-5.9,-2.2) {terrier};
  \node[w] (bl) at (-3.6,-2.2) {bloodhound};
  \node[w] (co) at (-1.0,-2.2) {cocker spaniel};
  \node[w] (t1) at (-8.7,-3.3) {Cesky terrier};
  \node[w] (t2) at (-5.9,-3.3) {Border terrier};
  \node[w] (t3) at (-2.5,-3.3) {Australian terrier};
  \foreach \c in {do,ca,ho,pi} \draw[e] (an.south) -- (\c.north);
  \foreach \c in {te,bl,co} \draw[e] (do.south) -- (\c.north);
  \foreach \c in {t1,t2,t3} \draw[e] (te.south) -- (\c.north);
\end{tikzpicture}
\captionof{figure}{A hyponymy (IS-A) hierarchy: a Border terrier is a terrier, a terrier
is a dog, a dog is an animal. Hyponymy is transitive.}
\end{center}

\subsubsection{Meronymy and holonymy}
\slidesrc{Semantics (long version): 18--20}

\begin{nbdef}{Meronymy and holonymy}
A concept $C_1$ is a \textbf{meronym} of a concept $C_2$ in a language $L$ if native
speakers of $L$ accept sentences like ``A $C_2$ has a $C_1$ (as a part)'' or ``A $C_1$ is
a part of a $C_2$''. The inverse relation is \textbf{holonymy}. It is the
\textbf{HAS-AS-PART} relation, organised in a part hierarchy (e.g.\ \emph{car} has
\emph{engine}, \emph{door}, \emph{wheel}, \emph{tyres}; \emph{engine} has
\emph{cylinder}, \emph{exhaust valve}, \ldots).
\end{nbdef}

Part-of is asymmetric and, \emph{with caution}, transitive. Transitivity fails when
different kinds of part-whole relation are chained:
\begin{itemize}[nosep]
  \item A musician has an arm. \hfill (component of a body)
  \item An orchestra has a musician. \hfill (member of a group)
  \item but: ?\,An orchestra has an arm.
\end{itemize}

\subsubsection{Homonymy}
\slidesrc{Semantics (long version): 21}

\begin{nbdef}{Homonyms, homophones, homographs}
\begin{itemize}[nosep]
  \item \textbf{Homonyms}: the same word form with multiple unrelated meanings.
  \item \textbf{Homophones}: words that \emph{sound} alike but are written differently.
  \item \textbf{Homographs}: words with the same spelling, possibly different
        pronunciation, and different meanings.
\end{itemize}
\end{nbdef}

\begin{center}
\footnotesize
\begin{tabularx}{\linewidth}{@{}*{3}{>{\raggedright\arraybackslash}X}@{}}
\toprule
\textbf{Homonyms} & \textbf{Homophones} & \textbf{Homographs} \\
\midrule
the \emph{spruce} tree / to \emph{spruce} up & \emph{addition} for maths / \emph{edition} of a book & \emph{desert} = abandon / area of land \\
\emph{suit} yourself / wore a \emph{suit} & I want \emph{to} go / I like it \emph{too} / one plus one is \emph{two} & \emph{bass} = fish / instrument \\
weigh on the \emph{scale} / \emph{scale} the wall & \emph{capitol} building / state \emph{capital} & \emph{close} = nearby / to shut \\
the price is \emph{fair} / go to the \emph{fair} & pick a \emph{flower} / bake with \emph{flour} & \emph{bow} = to bend down / ribbon \\
\bottomrule
\end{tabularx}
\end{center}

\begin{nbwhy}
Homographs and homonyms are a problem for every task that looks at written forms: a
search for \emph{bass} returns both fish and guitars (low precision, Section~\ref{sec:precision-recall}),
and a distributional vector for \emph{bass} mixes the contexts of all its meanings.
Choosing the right meaning in context is word sense disambiguation
(Section~\ref{sec:word-sense-disambiguation}).
\end{nbwhy}

\subsection{Thesauri and WordNet}
\label{sec:wordnet}
\slidesrc{Semantics (long version): 22--30}

\subsubsection{Structured lexicons and thesauri}
\slidesrc{Semantics (long version): 22--23}

A \textbf{thesaurus} (structured lexicon) is an alternative to the alphabetical
dictionary: a list of words grouped by \emph{meaning}, with a \textbf{hierarchical}
organisation, like the taxonomies of the natural sciences (children are ``types of'' the
father and inherit its properties). For ordinary words the hierarchy is given by hyponymy
and synonymy (Section~\ref{sec:lexical-relations}).
\begin{itemize}[nosep]
  \item It shows the structure of lexical knowledge, and is much used for technical
        terminology.
  \item It can be enriched with other relations: antonyms, part/member relations
        (\emph{has-as-part}, \emph{has-as-member}), not just \emph{is-a-type-of}.
  \item It can be explored in any direction (up, down, across), and obvious
        \textbf{distance metrics} (e.g.\ the number of links between two words) measure
        their similarity.
\end{itemize}

\subsubsection{WordNet}
\slidesrc{Semantics (long version): 24--30}

\textbf{WordNet} is a large lexical database of English, started in \textbf{1985} at
Princeton University by a group of psychologists and linguists (led by George Miller,
the author of the quotation in Section~\ref{sec:distributional-hypothesis}). Its
theoretical basis is psycholinguistics: what are the properties of the \emph{mental
lexicon}? Global organisation: \textbf{nouns} are organised as topical hierarchies,
\textbf{verbs} by entailment relations.

\begin{nbdef}{Synset}
In WordNet the basic unit is the \textbf{synset} (synonym set). A word meaning is
represented simply by listing the word forms that can express it. Synsets are
unambiguous designators: the two main senses of \emph{board} are
\begin{center}
\{board, plank, \ldots\} \quad (a piece of lumber) \qquad vs.\qquad \{board, committee, \ldots\}
\quad (a group of people assembled for some purpose).
\end{center}
\end{nbdef}

\begin{itemize}[nosep]
  \item Members of a synset are rarely true synonyms: WordNet does not try to capture
        subtle distinctions (connotation, collocation) among them.
  \item Synsets are often enough to \emph{differentiate} senses. If no suitable synonym
        exists, a short \textbf{gloss} is used, e.g.\ \{board, (a person's meals,
        provided regularly for money)\}. WordNet gives a gloss for every meaning, and
        often an example sentence.
\end{itemize}

\begin{nbexample}{the noun senses of \emph{board} in WordNet (excerpt)}
\small
\begin{itemize}[nosep]
  \item board (a committee having supervisory powers) ``\emph{the board has seven
        members}'';
  \item board, plank (a stout length of sawn timber; made in a wide variety of sizes
        and used for many purposes);
  \item board (a flat piece of material designed for a special purpose) ``\emph{he
        nailed boards across the windows}'';
  \item board, table (food or meals in general) ``\emph{she sets a fine table}'';
  \item display panel, display board, board (a vertical surface on which information
        can be displayed to public view);
  \item dining table, board (a table at which meals are served);
  \item control panel, instrument panel, control board, board, panel (electrical device
        consisting of a flat insulated surface that contains switches and dials).
\end{itemize}
For the first sense WordNet also lists its relations: its members (\emph{director},
\emph{board member}: meronyms), its hypernym (\emph{committee}), its sister terms.
\end{nbexample}

\paragraph{Size.} WordNet 3.0:
\begin{center}
\footnotesize
\setlength{\tabcolsep}{4pt}
\begin{tabular}{@{}l rrr @{\hspace{0.6cm}} rrr@{}}
\toprule
 & \textbf{Words} & \textbf{Synsets} & \textbf{Word--sense} & \textbf{Monos.} & \textbf{Polys.} & \textbf{Polys.} \\
 &  &  & \textbf{pairs} & \textbf{words} & \textbf{words} & \textbf{senses} \\
\midrule
Noun      & 117\,798 & 82\,115  & 146\,312 & 101\,863 & 15\,935 & 44\,449 \\
Verb      & 11\,529  & 13\,767  & 25\,047  & 6\,277   & 5\,252  & 18\,770 \\
Adjective & 21\,479  & 18\,156  & 30\,002  & 16\,503  & 4\,976  & 14\,399 \\
Adverb    & 4\,481   & 3\,621   & 5\,580   & 3\,748   & 733     & 1\,832 \\
\midrule
Total     & 155\,287 & 117\,659 & 206\,941 & 128\,391 & 26\,896 & 79\,450 \\
\bottomrule
\end{tabular}
\end{center}
Most words have a single sense (128\,391 of 155\,287), but the polysemous ones are the
frequent ones: verbs in particular have on average $18\,770/5\,252 \approx 3.6$ senses
each.

\begin{center}
\small
\begin{tabularx}{\linewidth}{@{}l >{\raggedright\arraybackslash}X@{}}
\toprule
\textbf{POS} & \textbf{Lexical relations in WordNet} \\
\midrule
Nouns & synonym $\sim$ antonym (opposite of); hypernym (is a kind of) $\sim$ hyponym (for
  example); coordinate (sister) terms: share the same hypernym; holonym (is part of)
  $\sim$ meronym (has as part) \\
Verbs & synonym $\sim$ antonym; hypernym $\sim$ troponym (\emph{lisp} -- \emph{talk});
  entailment (\emph{snore} -- \emph{sleep}: if you snore you sleep); coordinate terms \\
Adjectives, adverbs & the above, plus related nouns, verb participles, derivational
  information \\
\bottomrule
\end{tabularx}
\end{center}

WordNet is available in NLTK (Section~\ref{sec:word-segmentation}) and online
(``WordNet Search''); for Italian and many other languages, multilingual resources such as
BabelNet extend it.

\subsection{Word Sense Disambiguation}
\label{sec:word-sense-disambiguation}
\slidesrc{Semantics (long version): 31--39}

\begin{nbdef}{Word sense disambiguation (WSD)}
\textbf{Word sense disambiguation} is the task of identifying which sense (meaning) of a
word is used in a sentence, when the word has several meanings. It is an open problem of
NLP and ontology engineering.
\end{nbdef}

Consider the senses of \emph{bass} (WordNet lists eight noun senses: the lowest part of
the musical range, the lowest part in polyphonic music, an adult male singer with the
lowest voice, the lean flesh of a saltwater fish, any of various North American
freshwater fish, the member with the lowest range of a family of instruments, \ldots) and
the two sentences
\begin{itemize}[nosep]
  \item I went fishing for some sea bass.
  \item The bass line of the song is too weak.
\end{itemize}
To a human it is obvious that the first sentence uses \emph{bass} (fish) and the second
\emph{bass} (sound/instrument). Replicating this ability with an algorithm is difficult,
as shown even by the implicit equivocation between ``bass (sound)'' and ``bass (musical
instrument)'': the senses themselves are not always clearly separated.

\subsubsection{The Lesk algorithm}
\slidesrc{Semantics (long version): 35--37}

The \textbf{Lesk algorithm} (Michael Lesk, 1986, ``Automatic sense disambiguation using
machine readable dictionaries: how to tell a pine cone from an ice cream cone'') is the
classical WSD algorithm. It assumes that words in a given ``neighbourhood'' (a section of
text) tend to share a common topic.

\begin{nbdef}{Simplified Lesk}
For every sense of the word to disambiguate, count the words that occur both in the
\textbf{neighbourhood} of the word and in the \textbf{dictionary definition} (gloss) of
that sense. Choose the sense with the largest count.
\end{nbdef}

\begin{nbexample}{pine cone}
In the original version the glosses of \emph{two} ambiguous neighbouring words are
compared with each other. Context: ``pine cone''. Dictionary definitions:
\begin{center}
\small
\begin{tabular}{@{}l l@{}}
\toprule
\textbf{pine} 1 & kinds of \underline{evergreen} \underline{tree} with needle-shaped leaves \\
\textbf{pine} 2 & waste away through sorrow or illness \\
\midrule
\textbf{cone} 1 & solid body which narrows to a point \\
\textbf{cone} 2 & something of this shape whether solid or hollow \\
\textbf{cone} 3 & fruit of certain \underline{evergreen} \underline{trees} \\
\bottomrule
\end{tabular}
\end{center}
The best intersection is pine 1 $\cap$ cone 3 $=$ \{evergreen, tree(s)\}, with 2 words:
\emph{pine} is the tree and \emph{cone} is its fruit. All other pairs share no content
word.
\end{nbexample}

\begin{nbwarning}[tree vs.\ trees]
The overlap of 2 counts \emph{tree} and \emph{trees} as the same word. A literal string
comparison would find only \emph{evergreen}. Lesk needs the normalisation of
Section~\ref{sec:stemming-lemmatization} (and stop-word removal,
Section~\ref{sec:stop-words-tfidf}, otherwise \emph{of} and \emph{which} would count as
overlaps too).
\end{nbwarning}

\begin{nbexample}{simplified Lesk on ``I went fishing for some sea bass''}
After removing stop words and stemming, the context is \{go, fish, sea\}. Glosses
(with the synset members) of three senses of \emph{bass}:
\begin{itemize}[nosep]
  \item \emph{sea bass, bass}: the lean flesh of a saltwater fish of the family
        Serranidae $\to$ overlap \{sea, fish\} $= 2$;
  \item \emph{freshwater bass, bass}: any of various North American freshwater fish with
        lean flesh $\to$ overlap \{fish\} $= 1$;
  \item \emph{bass}: the lowest part of the musical range $\to$ overlap $0$.
\end{itemize}
Simplified Lesk chooses the sea bass. Note that it found \emph{sea} only thanks to the
synset member \emph{sea bass}: \emph{saltwater} and \emph{sea} do not match as strings.
\end{nbexample}

\subsubsection{Extensions and limitations}
\slidesrc{Semantics (long version): 38--39}

\begin{itemize}
  \item \textbf{Adapted Lesk} (Banerjee and Pedersen, 2002): uses WordNet, comparing
        also the glosses of related synsets (hypernyms, hyponyms, \ldots), which gives
        more words to overlap.
  \item \textbf{Limitations}: Lesk is very sensitive to the exact wording of
        definitions: the absence of a single word can change the result radically. And it
        only looks at overlaps among the glosses of the senses considered, which are
        short and do not provide enough vocabulary for fine-grained sense distinctions.
  \item Many later variants use other resources (thesauri, synonym dictionaries,
        morphological and syntactic models), and today distributional/contextual
        representations (Section~\ref{sec:distributional-semantics}) are the standard
        approach.
\end{itemize}

\subsection{Named Entity Recognition}
\label{sec:named-entity-recognition}
\slidesrc{Semantics (long version): 40--43}

\begin{nbdef}{Named entity recognition (NER)}
\textbf{Named entity recognition} locates and classifies the \textbf{named entities} of a
text into predefined categories: names of persons, organisations, locations, expressions
of times, quantities, monetary values, percentages, etc.
\end{nbdef}

Most NER systems take an unannotated block of text and produce an annotated one:
\begin{center}
Jim bought 300 shares of Acme Corp.\ in 2006.\\[3pt]
$\downarrow$\\[3pt]
[Jim]$_{\textsc{person}}$ bought 300 shares of [Acme Corp.]$_{\textsc{organization}}$ in
[2006]$_{\textsc{time}}$.
\end{center}

Online demos: Stanford NER (\url{http://nlp.stanford.edu:8080/ner/process}), spaCy NER;
for Italian (not only NER): \url{http://www.italianlp.it/},
\url{http://parli.di.unito.it/link_it.html}.

\paragraph{Approaches.}
\begin{itemize}[nosep]
  \item \textbf{Grammar-based} (hand-crafted rules by linguists): typically better
        \textbf{precision}, but lower \textbf{recall} and months of work by experienced
        computational linguists.
  \item \textbf{Statistical} (machine learning): need a large amount of manually
        annotated training data.
  \item \textbf{Semi-supervised} approaches avoid part of the annotation effort.
\end{itemize}

\begin{nbintuition}[why NER is not a dictionary lookup]
A list of names is not enough: \emph{Washington} can be a person, a city, a state or a
government (``Washington announced\ldots''); \emph{Apple} can be a fruit or a company;
new names appear every day. NER needs context, like POS tagging and WSD: in fact it is
often solved with the same sequence models (HMMs, MEMMs and their successors,
Section~\ref{sec:hmm-pos-tagging}), with tags such as B-PER, I-PER, O marking the
beginning and inside of each entity.
\end{nbintuition}

\begin{nbsummary}[lexical semantics, WSD, NER]
\begin{itemize}[nosep]
  \item Language = lexicon + grammar. Lexical relations: \textbf{synonymy} (substitution
        without changing truth, usually only in context), \textbf{antonymy} (between word
        forms), \textbf{hyponymy/hypernymy} (IS-A; troponymy for verbs),
        \textbf{meronymy/holonymy} (HAS-AS-PART; transitive only with caution),
        homonyms/homophones/homographs.
  \item Thesauri organise words by meaning in hierarchies. \textbf{WordNet}: synsets +
        glosses + relations; nouns in IS-A hierarchies, verbs by entailment.
  \item \textbf{WSD}: choose the sense in context. \textbf{Lesk}: the sense whose gloss
        overlaps most with the context (or with the glosses of neighbouring words);
        sensitive to wording, needs normalisation.
  \item \textbf{NER}: find and classify names, times, quantities; rules (precision) vs.\
        statistical (needs annotated data) vs.\ semi-supervised.
\end{itemize}
\end{nbsummary}

